\chapter{LINEAR ALGEBRA AND THE SPECTRAL THEOREM}

The emphasis in beginning analysis courses is on the behavior of individual (real or complex valued) functions.  In a \emph{functional analysis}
course the focus is shifted to \emph{spaces} of such functions and certain classes of mappings between these spaces.  It turns out that a concise,
perhaps overly simplistic, but certainly not misleading, description of such material is \emph{infinite dimensional linear algebra}.  From an
analyst's point of view the greatest triumphs of linear algebra were, for the most part, theorems about operators on finite dimensional vector
spaces (principally $\R^n$ and $\C^n$).  However, the vector spaces of scalar-valued mappings defined on various domains are typically infinite
dimensional.  Whether one sees the process of generalizing the marvelous results of finite dimensional linear algebra to the infinite dimensional
realm as a depressingly unpleasant string of abstract complications or as a rich source of astonishing insights, fascinating new questions, and
suggestive hints towards applications in other fields depends entirely on one's orientation towards mathematics in general.

In light of the preceding remarks it may be helpful to begin our journey into this new terrain with a brief review of some of the classical
successes of linear algebra.  It is not unusual for students who successfully complete a beginning linear algebra to have at the end only the
vaguest idea of what the course was about.  Part of the blame for this may be laid on the many elementary texts which make an unholy conflation
of two quite different, if closely related subjects: the study of \emph{vector spaces} and the linear maps between them on the one hand, and
\emph{inner product spaces} and their linear maps on the other.  Vector spaces have none of the geometric/topological notions of \emph{distance}
or \emph{length} or \emph{perpendicularity} or \emph{open} set or \emph{angle} between vectors; there is only addition and scalar multiplication.
Inner product spaces are vector spaces endowed with these additional structures.  Let's look at them separately.


\section{Vector Spaces and the Decomposition of Diagonalizable Operators}
\begin{conv} In these notes all vector
 \index{conventions!vector spaces are complex or real}%
 \index{C@$\C$!field of complex numbers}%
 \index{R@$\R$!field of real numbers}%
spaces will be assumed to be vector spaces over the field~$\C$ of complex numbers (in which case it is called a \emph{complex vector space}) or the
field~$\R$ of real numbers (in which case it is a \emph{real vector space}).  No other fields will appear.  When $\K$ appears it can be taken
 \index{K@$\K$ (field of real or complex numbers)}%
to be either $\C$ or~$\R$. A
 \index{scalar}%
\df{scalar} is a member of $\K$; that is, either a \emph{complex number} or a \emph{real number}.
\end{conv}

\begin{defn}\label{00001} The triple $(V,+,M)$ is a
 \index{vector!space}%
 \index{space!vector}%
\df{vector space} over $\K$ if $(V,+)$ is an Abelian group and $M\colon \K \sto \Hom(V)$
is a unital ring homomorphism (where $\Hom(V)$ is the ring of group homomorphisms
on~$V$).
\end{defn}

To check on the meanings of any the terms in the preceding definition, take a look at the first three sections of the first chapter of my
linear algebra notes~\cite{Erdman:2010}.

\begin{exer} The definition of \emph{vector space} found in many elementary texts is
something like the following: a \emph{vector space over} $\K$ is a set $V$ together with operations
of addition and scalar multiplication which satisfy the following axioms:
 \begin{enumerate}
  \item[(1)] if $x$, $y \in V$, then $x + y \in V$;
  \item[(2)] $(x + y) + z = x + (y + z)$ for every
 \index{associative}%
$x$, $y$, $z \in V$ (associativity);
  \item[(3)] there exists $\vc 0 \in V$ such that $x + \vc 0 = x$ for every
 \index{identity!additive}%
$x \in V$ (existence of additive identity);
  \item[(4)] for every $x\in V$ there exists $-x \in V$ such that
 \index{additive!inverses}%
 \index{inverse!additive}%
$x + (-x) = \vc 0$ (existence of additive inverses);
  \item[(5)] $x + y = y + x$ for every $x$, $y \in V$
 \index{commutative}%
(commutativity);
  \item[(6)] if $\alpha \in \K$ and $x \in V$, then $\alpha x \in V$;
  \item[(7)] $\alpha (x + y) = \alpha x + \alpha y$ for every $\alpha \in \K$
and every $x$, $y \in V$;
  \item[(8)] $(\alpha + \beta)x = \alpha x + \beta x$ for every $\alpha$,
$\beta \in \K$ and every $x \in V$;
  \item[(9)] $(\alpha\beta)x = \alpha(\beta x)$ for every $\alpha$,
$\beta \in \K$ and every $x \in V$; and
  \item[(10)] $1\,x = x$ for every $x \in V$.
 \end{enumerate}
Verify that this definition is equivalent to the one given above in~\ref{00001}.
\end{exer}

\begin{defn} A subset $M$ of a vector space $V$ is a
 \index{subspace!vector}%
 \index{vector!subspace}%
\df{vector subspace} of $V$ if it is a vector space under the operations it inherits from~$V$.
\end{defn}

\begin{prop} A nonempty subset of $M$ of a vector space $V$ is a vector subspace of $V$ if and
only if it is closed under addition and scalar multiplication. (That is: if $\vc x$ and
$\vc y$ belong to $M$, so does $\vc x + \vc y$; and if $\vc x$ belongs to $M$ and $\alpha
\in \K$, then $\alpha \vc x$ belongs to~$M$.)
\end{prop}

\begin{defn} A vector $y$ is a
 \index{linear!combination}%
 \index{combination!linear}%
\df{linear combination} of vectors $x_1$, \dots, $x_n$ if there exist scalars $\alpha_1$,
\dots $\alpha_n$ such that $y = \sum_{k=1}^n \alpha_k x_k$. The linear combination $\sum_{k=1}^n \alpha_k x_k$ is
 \index{linear!combination!trivial}%
 \index{trivial!linear combination}%
\df{trivial} if all the coefficients $\alpha_1$, \dots $\alpha_n$ are zero.  If at least
one $\alpha_k$ is different from zero, the linear combination is \df{nontrivial}.
\end{defn}

\begin{defn} If $A$ is a nonempty subset of a vector space~$V$, then $\spn A$, the
 \index{span}%
\df{span} of $A$, is the set of all linear combinations of elements of~$A$.  This is also frequently called the
 \index{span!linear}%
 \index{linear!span}%
\df{linear span} of~$A$.
\end{defn}

\begin{notn}\label{bases111} Let $A$ be a set which spans a vector space~$V$.  For every
$x \in V$ there exists a finite set $S$ of vectors in $A$ and for each element $e$ in $A$
there exists a scalar $x_e$ such that $x = \sum_{e \in S} x_e e$.  If we agree to let $x_e = 0$
whenever $e \in B \setminus S$, we can just as well write $x = \sum_{e \in B} x_e e$.  Although
this notation may make it appear as if we are summing over an arbitrary, perhaps uncountable,
set, the fact of the matter is that all but finitely many of the terms are zero, so no
``convergence'' problems arise. Treat $\sum_{e \in B} x_e e$ as a finite sum.
Associativity and commutativity of addition in $V$ make the expression unambiguous.
\end{notn}

\begin{defn}\label{defn_lin_dep} A subset $A$ of a vector space is
 \index{linear!dependence}%
 \index{dependent}
\df{linearly dependent} if the zero vector $\vc 0$ can be written as a nontrivial linear
combination of elements of~$A$; that is, if there exist vectors $x_1, \dots, x_n \in A$
and scalars $\alpha_1, \dots, \alpha_n$, \textbf{not all zero,} such that $\sum_{k=1}^n
\alpha_kx_k = \vc 0$. A subset of a vector space is
 \index{linear!independence}%
 \index{independent}
\df{linearly independent} if it is not linearly dependent.
\end{defn}

Technically, it is a \emph{set} of vectors that is linearly dependent or independent.
Nevertheless, these terms are frequently used as if they were properties of the vectors
themselves.  For instance, if $S = \{x_1, \dots, x_n\}$ is a finite set of vectors in a
vector space, you may see the assertions ``the set $S$ is linearly independent'' and
``the vectors $x_1$, \dots $x_n$ are linearly independent'' used interchangeably.

\begin{defn}\label{defn_Hamel_basis} A set $B$ of vectors in a vector space $V$ is a
 \index{basis!Hamel}%
 \index{Hamel basis}%
\df{Hamel basis} for $V$ if it is linearly independent and spans~$B$.
\end{defn}

\begin{conv}\label{conv_1_1} The vector space containing a single vector, the zero vector, is the
 \index{convention!basis for the zero vector space}%
 \index{vector!space!trivial (or zero)}%
 \index{trivial!vector space}%
\df{trivial} (or
 \index{zero!vector space}%
\df{zero}) vector space. Associated with this space is a technical, if not very interesting, question.  Does it have a Hamel basis?
It is clear from the definition of \emph{linear dependence}~\ref{defn_lin_dep} that the empty set $\emptyset$ is linearly
independent.  And it is certainly a maximal linearly independent subset of the zero space (which condition, for nontrivial
spaces, is equivalent to being a linearly independent spanning set).  But it is hard to argue that the empty set spans anything.
So, simply as a matter of convention, we will say that the answer to the question is \emph{yes}.  The zero vector space has a Hamel
basis, and it is the empty set!
\end{conv}

In beginning linear algebra texts a linearly independent spanning set is called simply a \emph{basis}.  I use \emph{Hamel basis}
instead to distinguish it from \emph{orthonormal basis} and \emph{Schauder basis}, terms which refer to concepts which are of
more interest in functional analysis.

\begin{prop} Let $B$ be a Hamel basis for a nontrivial vector space~$V$.  Then every element in $V$ can be written in a unique way as a linear
combination of members of~$B$.
\end{prop}

\begin{proof}[\emph{Hint for proof}] Existence is clear.  Simplify the proof of uniqueness by making use of the notational
convention suggested in~\ref{bases111}.   \ns
\end{proof}

\begin{prop}\label{bases025} Let $A$ be a linearly independent subset of a vector
space~$V$. Then there exists a Hamel basis for~$V$ which contains~$A$.
\end{prop}

\begin{proof}[\emph{Hint for proof}] Show first that, for a nontrivial space, a linearly independent subset of $V$ is a basis for~$V$
if and only if it is a maximal linearly independent subset. Then order the family of linearly independent subsets of
$V$ which contain~$A$ by inclusion and apply \emph{Zorn's lemma}. (If you are not familiar with \emph{Zorn's lemma},
see section 1.7 of my linear algebra notes~\cite{Erdman:2010}.)  \ns
\end{proof}

\begin{cor} Every vector space has a Hamel basis.
\end{cor}

\begin{defn}\label{000082} Let $M$ and $N$ be subspaces of a vector space $V$. If
$M \cap N = \{\vc 0\}$ and $M + N = V$, then $V$ is the
 \index{internal!direct sum!in vector spaces}%
 \index{direct!sum!internal}%
 \index{direct!sum!of vector spaces}%
 \index{<binopdirectsum@$\oplus$ (vector space direct sum)}%
(\df{internal}) \df{direct sum} of $M$ and $N$. In this case we write
   \[ V = M \oplus N\,. \]
We say that $M$ and $N$ are
 \index{complement!of a vector subspace}%
 \index{vector!space!complement}%
\df{complementary vector subspaces} and that each is a (vector space) \df{complement} of the
other. The
 \index{codimension}%
\df{codimension} of the subspace $M$ is the dimension of its complement~$N$.
\end{defn}

\begin{prop}\label{0000823} Let $V$ be a vector space and suppose that $V = M \oplus N$.  Then for
every $v \in V$ there exist unique vectors $m \in M$ and $n \in N$ such that $v = m + n$.
\end{prop}

\begin{exam} Let $\fml C = \fml C[-1,1]$ be the vector space of all continuous real valued
functions on the interval $[-1,1]$. A function $f$ in $\fml C$ is \df{even} if $f(-x) =
f(x)$ for all $x \in [-1,1]$; it is \df{odd} if $f(-x) = -f(x)$ for all $x \in [-1,1]$.
Let $\fml C_o = \{f \in \fml C\colon  f \text{ is odd }\}$ and $\fml C_e = \{f \in \fml
C\colon  f \text{ is even }\}$. Then $\fml C = \fml C_o \oplus \fml C_e$.
\end{exam}

\begin{prop}\label{0000824} If $M$ is a subspace of a vector space $V$, then there exists a subspace $N$ of
$V$ such that $V = M \oplus N$.
\end{prop}

\begin{lem}\label{bases030} Let $V$ be a vector space with a nonempty finite basis  $\{e^1,
\dots,e^n\}$ and let $v = \sum_{k=1}^n \alpha_k e^k$ be a vector in~$V$.  If $p \in \N_n$
and $\alpha_p \ne 0$, then $\{e^1, \dots,e^{p-1},v,e^{p+1}, \dots, e^n\}$ is a basis
for~$V$.
\end{lem}

\begin{prop}\label{bases031} If some basis for a vector space $V$ contains $n$ elements,
then every linearly independent subset of $V$ with $n$ elements is also a basis.
\end{prop}

\begin{proof}[\emph{Hint for proof}] Suppose $\{e^1,\dots,e^n\}$ is a basis for~$V$ and
$\{v^1, \dots,v^n\}$ is linearly independent in~$V$. Start by using lemma~\ref{bases030}
to show that (after perhaps renumbering the $e^k$'s) the set $\{v^1, e^2, \dots, e^n\}$
is a basis for~$V$.  \ns
\end{proof}

\begin{cor}\label{bases032} If a vector space $V$ has a finite basis $B$, then every basis
for $V$ is finite and contains the same number of elements as~$B$.
\end{cor}

\begin{defn} A vector space is
 \index{finite dimensional}%
\df{finite dimensional} if it has a finite basis and the
 \index{dimension!of a vector space}%
 \index{vector!space!dimension of a}%
 \index{dim@$\dim V$ (dimension of a vector space $V$)}%
\df{dimension} of the space is the number of elements in this (hence any) basis for the
space. The dimension of a finite dimensional vector space $V$ is denoted by $\dim V$. By the
convention made above~\ref{conv_1_1}, the zero vector space has dimension zero.  If
the space does not have a finite basis, it is
 \index{infinite dimensional}%
\df{infinite dimensional}.
\end{defn}

\begin{defn}\label{defn_op_vsp} A function $T\colon V \sto W$ between vector spaces (over the same field) is
 \index{linear!map}%
\df{linear} if $T(u + v) = Tu + Tv$ for all $u$, $v \in V$ and $T(\alpha v) = \alpha Tv$
for all $\alpha \in \K$ and $v \in V$. Linear functions are frequently called
\emph{linear transformations} or \emph{linear maps}. When $V = W$ we say that the linear map $T$ is an
 \index{operator!on a vector space}%
\df{operator} on~$V$. Depending on context we denote the identity operator $x \mapsto x$ on $V$
 \index{identity@$\id V$ or $I_V$ or $I$ (identity operator)}%
by $\id V$ or $I_V$ or just~$I$. Recall that if $T\colon V \sto W$ is a linear map, then the
 \index{kernel!of a linear map}%
 \index{<ker@$\ker T$ (kernel of~$T$)}%
\df{kernel} of $T$, denoted by $\ker T$, is $T^{\gets}(\{0\}) := \{x \in V\colon Tx = \vc
0\}$. Also, the
 \index{range!of a linear map}%
 \index{<ran@$\ran T$ (range of~$T$)}%
\df{range} of $T$, denoted by $\ran T$, is $T^{\sto}(V) := \{Tx\colon x \in V\}$.
\end{defn}

\begin{prop} A linear map $T\colon V \sto W$ is injective if and only if its kernel is zero.
\end{prop}

\begin{prop} Let $T\colon V \sto W$ be an injective linear map between vector spaces. If $A$ is a linearly independent
subset of~$V$, then $T^\sto(A)$ is a linearly independent subset of~$W$.
\end{prop}

\begin{prop} Let $T\colon V \sto W$ be a linear map between vector spaces and $A \subseteq V$.
Then $T^\sto(\spn A) = \spn T^\sto(A)$.
\end{prop}

\begin{prop} Let $T\colon V \sto W$ be an injective linear map between vector spaces. If $B$ is a basis
for a subspace $U$ of~$V$, then $T^\sto(B)$ is a basis for~$T^\sto(U)$.
\end{prop}

\begin{prop} Two finite dimensional vector spaces are isomorphic if and only if they have the same dimension.
\end{prop}

\begin{defn} Let $T$ be a linear map between vector spaces. Then $\textrm{rank }T$, the
 \index{rank}
\df{rank} of $T$ is the dimension of the range of ~$T$, and $\textrm{nullity }T$, the
 \index{nullity}%
\df{nullity} of $T$ is the dimension of the kernel of~$T$.
\end{defn}

\begin{prop} Let $T\colon V \sto W$ be a linear map between vector spaces. If $V$ is finite dimensional, then
  \[ \textrm{rank }T + \textrm{nullity }T = \dim V. \]
\end{prop}

\begin{defn} A linear map $T\colon V \sto W$ between vector spaces is
 \index{invertible!linear map}%
\df{invertible} (or is an
 \index{isomorphism!of vector spaces}%
\df{isomorphism}) if there exists a linear map $T^{-1}\colon W \sto V$ such that $T^{-1}T
= \id V$ and $TT^{-1} = \id W$.  Two vector spaces $V$ and $W$ are
 \index{isomorphic!vectorspaces}%
\df{isomorphic} if there exists an isomorphism from $V$ to~$W$.
\end{defn}

\begin{prop}\label{inv_lin_maps001} A linear map $T\colon V \sto W$ between vector spaces
is invertible if and only if has both a left inverse and a right inverse.
\end{prop}

\begin{prop}\label{inv_lin_maps004} A linear map between vector spaces is invertible if
and only if it is bijective.
\end{prop}

\begin{prop} For an operator $T$ on a finite dimensional vector space the following are equivalent:
 \begin{enumerate}
    \item $T$ is an isomorphism;
    \item $T$ is injective; and
    \item $T$ is surjective.
 \end{enumerate}
\end{prop}

 \begin{defn}\label{defn_similar_vs_ops} Two operators $R$ and $T$ on a vector space $V$ are
 \index{similar!operators}%
\df{similar} if there exists an invertible operator $S$ on $V$ such that $R = STS^{-1}$.
\end{defn}

\begin{prop}\label{sim004thm} If $V$ is a vector space, then similarity is an equivalence
relation on~$\ofml L(V)$.
\end{prop}

Let $V$ and $W$ be finite dimensional vector spaces with ordered bases. Suppose that $V$ is $n$-dimensional with ordered basis
$\{e^1,e^2, \dots,e^n\}$ and $W$ is $m$-dimensional. Recall from beginning linear algebra that if $T\colon V \sto W$ is linear,
then its
 \index{matrix!representation}%
 \index{representation!matrix}%
\emph{matrix representation} $[T]$ (taken with respect to the ordered bases in $V$ and~$W$) is the $m \times n$-matrix $[T]$
whose $k^{\text{th}}$ column ($1 \le k \le n$) is the column vector $Te^k$ in~$W$.  The point here is that the action of $T$ on
a vector $x$ can be \emph{represented} as multiplication of $x$ by the matrix $[T]$; that is,
   \[ Tx = [T]x. \]
Perhaps this equation requires a little interpretation. The left side is the function $T$ evaluated at $x$, the
result of this evaluation being thought of as a column vector in~$W$; the right side is an $m \times n$ matrix multiplied
by an $n \times 1$ matrix (that is, a column vector).  So the asserted equality is of two $m \times 1$ matrices (column vectors).

\begin{defn} Let $V$ be a finite dimensional vector space and $B = \{e^1, \dots, e^n\}$ be a basis for~$V$.  An operator $T$ on $V$ is
 \index{diagonal!operator}%
 \index{operator!diagonal}%
\df{diagonal} if there exist scalars $\alpha_1, \dots, \alpha_n$ such that $Te^k =
\alpha_ke^k$ for each $k \in \N_n$. Equivalently, $T$ is diagonal if its matrix
representation $[T] = [T_{ij}]$ has the property that $T_{ij} = 0$ whenever $i \ne j$.
\end{defn}

Asking whether a particular operator on some finite dimensional vector space is diagonal
is, strictly speaking, nonsense. As defined, the operator property of being diagonal is
definitely \emph{not} a vector space concept. It makes sense only for a vector space
\emph{for which a basis has been specified}.  This important, if obvious, fact seems to
go unnoticed in many beginning linear algebra courses, due, I suppose, to a rather obsessive
fixation on $\R^n$ in such courses. Here is the relevant \emph{vector space} property.

\begin{defn} An operator $T$ on a finite dimensional vector space $V$ is
 \index{diagonalizable!operator}%
 \index{operator!diagonalizable}%
\df{diagonalizable} if there exists a basis for $V$ with respect to which $T$ is
diagonal. Equivalently, an operator on a finite dimensional vector space \emph{with
basis} is diagonalizable if it is similar to a diagonal operator.
\end{defn}

\begin{defn} Let $V$ be a vector space and suppose that $V = M \oplus N$.  We know
from~\ref{0000823} that for each $\vc v \in V$ there exist unique vectors $\vc m \in M$
and $\vc n \in N$ such that $\vc v = \vc m + \vc n$.  Define a function $\sbsb E{MN}\colon V \sto V$ by $\sbsb E{MN}v = n$.
The function $\sbsb E{MN}$ is the
 \index{projection!in a vector space}%
 \index{E@$\sbsb E{MN}$ (projection along $M$ onto $N$)}%
\df{projection of} $V$ \df{along} $M$ \df{onto}~$N$.  (Frequently we write $E$ for $\sbsb E{MN}$. But keep in mind that
$E$ depends on both $M$ and~$N$.)
\end{defn}

\begin{prop} Let $V$ be a vector space and suppose that $V = M \oplus N$. If $E$ is the
projection of $V$ along $M$ onto $N$, then
 \begin{enumerate}[label=\rm{(\alph*)}]
  \item $E$ is linear;
 \index{idempotent}%
  \item $E^2 = E$ \quad (that is, $E$ is \df{idempotent});
  \item $\ran E = N$; and
  \item $\ker E = M$.
 \end{enumerate}
\end{prop}

\begin{prop} Let $V$ be a vector space and suppose that $E\colon V \sto V$ is a function which
satisfies
 \begin{enumerate}[label=\rm{(\alph*)}]
  \item $E$ is linear, and
  \item $E^2 = E$.
\end{enumerate}
Then
 \[V = \ker E \oplus \ran E\]
and $E$ is the projection of $V$ along $\ker E$ onto~$\ran E$.
\end{prop}

It is important to note that an obvious consequence of the last two propositions is that
a function $T\colon V \sto V$ from a finite dimensional vector space into itself is a
projection if and only if it is linear and idempotent.







\begin{prop} Let $T\colon V \sto W$ be a linear transformation between vector spaces. Then
  \begin{enumerate}[label=\rm{(\alph*)}]
    \item $T$ has a left inverse if and only if it is injective; and
    \item $T$ has a right inverse if and only if it is surjective.
  \end{enumerate}
\end{prop}














\begin{prop} Let $V$ be a vector space and suppose that $V = M \oplus N$.  If $E$ is the
projection of $V$ along $M$ onto $N$, then $I - E$ is the projection of $V$ along $N$
onto~$M$.
\end{prop}

As we have just seen, if $E$ is a projection on a vector space $V$, then the identity
operator on $V$ can be written as the sum of two projections $E$ and $I - E$ whose
corresponding ranges form a direct sum decomposition of the space $V = \ran E \oplus \ran
(I - E)$.  We can generalize this to more than two projections.

\begin{defn} Suppose that on a vector space $V$ there exist projection operators $E_1$,
\dots, $E_n$ such that
\begin{enumerate}
    \item $I_V = E_1 + E_2 + \dots + E_n$ and
    \item $E_iE_j = 0$ whenever $i \ne j$.
\end{enumerate}
Then we say that $I_V = E_1 + E_2 + \dots + E_n$ is a
 \index{resolution of the identity!in vector spaces}%
\df{resolution of the identity}.
\end{defn}

\begin{prop} If $I_V = E_1 + E_2 + \dots + E_n$ is a resolution of the identity on a vector
space $V$, then $V = \bigoplus_{k=1}^n \ran E_k$.
\end{prop}

\begin{exam} Let $P$ be the plane in $\R^3$ whose equation is $x - z = 0$ and $L$ be the line
whose equations are $y = 0$ and $x = -z$.  Let $E$ be the projection of $\R^3$ along~$L$
onto $P$ and $F$ be the projection of $\R^3$ along~$P$ onto~$L$.  Then
 \[[E] = \begin{bmatrix}
              \frac12  & 0 & \frac12 \\
                  0    & 1 &   0     \\
              \frac12  & 0 & \frac12
         \end{bmatrix} \qquad \text{ and } \qquad
   [F] = \begin{bmatrix}
              \frac12  & 0 & -\frac12 \\
                  0    & 0 &   0     \\
             -\frac12  & 0 &  \frac12
         \end{bmatrix}\,. \]
\end{exam}

\begin{defn} A complex number $\lambda$ is an
 \index{eigenvalue}%
\df{eigenvalue} of an operator $T$ on a vector space $V$ if $\ker(T - \lambda I_V)$
contains a nonzero vector.  Any such vector is an
 \index{eigenvector}%
\df{eigenvector} of $T$ associated with $\lambda$ and $\ker(T - \lambda I_V)$ is the
 \index{eigenspace}
\df{eigenspace} of $T$ associated with~$\lambda$.  The set of all eigenvalues of the
operator $T$ is its
 \index{point spectrum}%
 \index{spectrum!point}%
\df{point spectrum} and is denoted by~$\sigma_p(T)$.
\end{defn}

If $M$ is an $n \times n$ matrix, then $\det(M - \lambda I_n)$ (where $I_n$ is the $n
\times n$ identity matrix) is a polynomial in $\lambda$ of degree~$n$. This is the
\df{characteristic polynomial} of~$M$. A standard way of computing the eigenvalues of an
operator $T$ on a finite dimensional vector space is to find the zeros of the
characteristic polynomial of its matrix representation. It is an easy consequence of the
multiplicative property of the determinant function that the characteristic polynomial of
an operator $T$ on a vector space $V$ is independent of the basis chosen for $V$ and
hence of the particular matrix representation of~$T$ that is used.

\vskip 5 pt

\begin{exam} The eigenvalues of the operator on (the real vector space) $\R^3$ whose matrix
representation is  $ \begin{bmatrix} 0  &  0  &  2 \\
                                     0  &  2  &  0 \\
                                     2  &  0  &  0
                      \end{bmatrix}$ are $-2$ and $+2$, the latter having (both algebraic and
geometric) multiplicity~$2$. The eigenspace associated with the negative eigenvalue is
$\spn\{(1,0,-1)\}$ and the eigenspace associated with the positive eigenvalue is
$\spn\{(1,0,1), (0,1,0)\}$.
\end{exam}

\vskip 5 pt

\begin{prop} Every operator on a complex finite dimensional vector space has an eigenvalue.
\end{prop}

\begin{exam} The preceding proposition does not hold for real finite dimensional spaces.
\end{exam}

\begin{proof}[\emph{Hint for proof}] Consider rotations of the plane.  \ns
\end{proof}

The central fact asserted by the finite dimensional vector space version of the
\emph{spectral theorem} is that every diagonalizable operator on such a space can be
written as a linear combination of projection operators where the coefficients of the
linear combination are the eigenvalues of the operator and the ranges of the projections
are the corresponding eigenspaces.  Thus if $T$ is a diagonalizable operator on a finite
dimensional vector space $V$, then $V$ has a basis consisting of eigenvectors of~$T$.

Here is a formal statement of the theorem.

\begin{thm}[Spectral Theorem: finite dimensional vector space version] Suppose that $T$ is a
 \index{spectral!theorem!for diagonalizable operators on vector spaces}%
diagonalizable operator on a finite dimensional vector space~$V$.  Let $\lambda_1, \dots, \lambda_n$ be
the (distinct) eigenvalues of~$T$. Then there exists a resolution of the identity $I_V =
E_1 + \dots + E_n$, where for each $k$ the range of the projection $E_k$ is the
eigenspace associated with~$\lambda_k$, and furthermore
     \[ T = \lambda_1E_1 +\dots + \lambda_nE_n\,. \]
\end{thm}

\begin{proof} A proof of this theorem can be found in \cite{HoffmanK:1971} on page~212. \ns
\end{proof}

\begin{exam}  Let $T$ be the operator on (the real vector space) $\R^2$ whose matrix representation
is $\begin{bmatrix} -7 &   8 \\ -16 &  17 \end{bmatrix}$.

\vskip 3 pt

 \begin{enumerate}
  \item The characteristic polynomial for~$T$ is
                     $\sbsb cT(\lambda) = \lambda^2 - 10\lambda + 9$.

\vskip 3 pt

  \item The eigenspace $M_1$ associated with the eigenvalue $1$ is $\spn \{(1,1)\}$.

\vskip 3 pt

  \item The eigenspace $M_2$ associated with the eigenvalue $9$ is $\spn \{(1,2)\}$.

\vskip 3 pt

  \item We can write $T$ as a linear combination of projection operators.  In particular,
   \[ T = 1\cdot E_1 + 9\cdot E_2 \text{ where } [E_1] = \begin{bmatrix} 2 & -1 \\ 2 & -1 \end{bmatrix}
\text{ and } [E_2] = \begin{bmatrix} -1 & 1 \\ -2 & 2 \end{bmatrix}\,. \]

\vskip 3 pt

  \item Notice that the sum of $[E_1]$ and $[E_2]$ is the identity matrix and that their
product is the zero matrix.

\vskip 3 pt

  \item The matrix  $S = \begin{bmatrix} 1 &  1  \\ 1 &  2 \end{bmatrix}$  diagonalizes
$[T]$. That is, $S^{-1}\,[T]\,S = \begin{bmatrix} 1 &  0  \\ 0 &  9 \end{bmatrix}$.


\vskip 3 pt


  \item A matrix representing $\sqrt T$ is $\begin{bmatrix} -1 &  2  \\ -4 &  5 \end{bmatrix}$.
 \end{enumerate}
\end{exam}

\begin{exer} Let $T$ be the operator on $\R^3$ whose matrix representation is
    \smash[b]{$\begin{bmatrix}
               2  &  0  &  0 \\
              -1  &  3  &  2 \\
               1  & -1  &  0 \end{bmatrix}$.}
  \begin{enumerate}
    \item Find the characteristic and minimal polynomials of~$T$.
    \item What can be concluded from the form of the minimal polynomial?
    \item Find a matrix which diagonalizes~$T$. What is the diagonal form of $T$ produced by this matrix?
    \item Find (the matrix representation of)~$\sqrt T$.
  \end{enumerate}
\end{exer}

\begin{defn} A vector space operator $T$ is
 \index{nilpotent}%
\df{nilpotent} if $T^n = \vc 0$ for some $n \in \N$.
\end{defn}

An operator on a finite dimensional vector space need not be diagonalizable.  If it is not,
how close to diagonalizable is it?  Here is one answer.

\begin{thm} Let $T$ be an operator on a finite dimensional vector space~$V$.  Suppose that the minimal
polynomial for $T$ factors completely into linear factors
   \[ m_T(x) = (x - \lambda_1)^{r_1} \dots (x - \lambda_k)^{r_k} \]
where $\lambda_1, \dots \lambda_k$ are the (distinct) eigenvalues of~$T$. For each $j$ let
$W_j = \ker(T - \lambda_jI)^{r_j}$ and $E_j$ be the projection of $V$ onto $W_j$ along
$W_1 + \dots + W_{j-1} + W_{j+1} + \dots + W_k$. Then
   \[ V = W_1 \oplus W_2 \oplus \dots \oplus W_k, \]
each $W_j$ is invariant under $T$, and $I = E_1 + \dots + E_k$. Furthermore, the operator
   \[ D = \lambda_1E_1 + \dots + \lambda_kE_k \]
is diagonalizable, the operator
   \[ N = T - D \]
is nilpotent, and $N$ commutes with~$D$.
\end{thm}

\begin{proof} See~\cite{HoffmanK:1971}, pages 222--223. \ns
\end{proof}

\begin{defn} Since, in the preceding theorem, $T = D + N$ where $D$ is diagonalizable and $N$ is nilpotent,
we say that $D$ is the
 \index{diagonalizable!part}%
 \index{part!diagonalizable}%
\df{diagonalizable part} of $T$ and $N$ is the
 \index{nilpotent!part}%
 \index{part!nilpotent}%
\df{nilpotent part} of~$T$.
\end{defn}

\begin{exer} Let $T$ be the operator on $\R^3$ whose matrix representation is
 $\begin{bmatrix}
        3  &  1  & -1 \\
        2  &  2  & -1 \\
        2  &  2  &  0
  \end{bmatrix}$.

 \begin{enumerate}
  \item Find the characteristic and minimal polynomials of~$T$.
  \item Find the eigenspaces of~$T$.
  \item Find the diagonalizable part $D$ and nilpotent part $N$ of~$T$.
  \item Find a matrix which diagonalizes~$D$. What is the diagonal form
of $D$ produced by this matrix?
  \item Show that $D$ commutes with~$N$.
 \end{enumerate}
\end{exer}














\section{Normal Operators on an Inner Product Space}
\begin{defn} Let $V$ be a vector space.  A function $s$ which associates to each pair of vectors
$x$ and $y$ in $V$ a scalar
 \index{s@$s(x,y)$ (semi-inner product)}%
 \index{<s@$s(x,y)$ (semi-inner product)}%
$s(x,y)$ is an
 \index{semi-inner product}%
 \index{semi-inner product!space}%
 \index{product!semi-inner}%
 \index{space!semi-inner product}%
\df{semi-inner product} on $V$ provided that for every $x$, $y$, $z \in V$ and $\alpha \in \K$ the following four conditions are satisfied:
 \begin{enumerate}
   \item $s(x + y , z) = s(x,z) + s(y,z)$;
   \item $s(\alpha x,y) = \alpha s(x,y)$;
   \item $s(x,y) = \conj{s(y,x)}$; and
   \item] $s(x,x) \ge 0$.
 \end{enumerate}
 If, in addition, the function $s$ satisfies
  \begin{enumerate}
    \item For every nonzero $x$ in $V$ we have $s(x,x) > 0$.
  \end{enumerate}
then $s$ is an
 \index{inner product}%
 \index{inner product!space}%
 \index{product!inner}%
 \index{space!inner product}%
\df{inner product} on~$V$.  We will usually write the inner product of two vectors $x$ and $y$ as
 \index{<bracpnt@$\langle x,y \rangle$ (inner product)}%
$\langle x,y \rangle$ rather than $s(x,y)$.

\noindent The overline in (c) denotes complex conjugation (so is redundant in the case of a real vector space).
Conditions (a) and (b) show that a semi-inner product is linear in its first variable.  Conditions (a) and (b) of
proposition~\ref{00015001} say that a complex inner product is
 \index{conjugate!linear}%
 \index{linear!conjugate}%
\df{conjugate linear} in its second variable.  When a scalar valued function of two variables on a complex semi-inner product
space is linear in one variable and conjugate linear in the other, it is often called
 \index{sesquilinear}%
 \index{linear!sesqui-}%
\df{sesquilinear}.  (The prefix ``sesqui-'' means ``one and a half''.)  We will also use the term \emph{sesquilinear} for a
bilinear function on a real semi-inner product space.  Taken together conditions (a)--(d) say that the inner product is a
\emph{positive definite, conjugate symmetric, sesquilinear form}.
\end{defn}

\begin{notn}\label{norm_notn} If $V$ is a vector space which has been equipped with an
semi-inner product and $x \in V$ we introduce the abbreviation
  \[ \norm x := \sqrt{s(x,x)} \]
which is read \emph{the norm of $x$} or \emph{the length of~$x$}.  (This somewhat
optimistic terminology is justified, at least when $s$ is an inner product, in proposition~\ref{00015025} below.)
\end{notn}

\begin{prop}\label{00015001} If $x$, $y$, and $z$ are vectors in a space with semi-inner product $s$ and $\alpha \in \K$, then
 \begin{enumerate}[label=\rm{(\alph*)}]
   \item $s(x,y + z)  = s(x,y) + s(x,z)$,
   \item $s(x,\alpha y) = \conj\alpha s(x,y)$, and
   \item $s(x,x) = 0$ if and only if $x = \vc 0$.
 \end{enumerate}
\end{prop}

\begin{exam}\label{000150012} For vectors $x = (x_1,x_2,\dots,x_n)$ and $y = (y_1,y_2,\dots,y_n)$
belonging to $\K^n$ define
   \[ \langle x,y \rangle = \sum_{k=1}^n x_k \conj{y_k}\,. \]
Then
 \index{K@$\K^n$!as an inner product space}%
 \index{inner product!space!$\K^n$ as a}%
$\K^n$ is an inner product space.
\end{exam}

\begin{exam}\label{000150013} Let $l_2 = l_2(\N)$ be the set of all square summable sequences of
complex numbers. (A sequence $x = (x_k)_{k=1}^\infty$ is
 \index{square summable}%
 \index{summable!square}%
 \index{l@$l_2 = l_2(\N)$!square summable sequences}%
 \index{l@$l_2 = l_2(\N)$!as an inner product space}%
 \index{inner product!space!$l_2$ as a}%
\df{square summable} if $\sum_{k=1}^\infty \abs{x_k}^2 < \infty$.) (The vector space operations are defined
pointwise.) For vectors $x = (x_1,x_2,\dots)$ and $y = (y_1,y_2,\dots)$ belonging to $l_2$ define
   \[ \langle x,y \rangle = \sum_{k=1}^\infty x_k \conj{y_k}\,. \]
Then $l_2$ is an inner product space.  (It must be shown, among other things, that the
series in the preceding definition actually converges.)  A similar space is $l_2(\Z)$ the set of all square summable
 \index{l@$l_2(\Z)$!square summable bilateral sequences}%
 \index{l@$l_2(\Z)$!as an inner product space}%
 \index{inner product!space!$l_2(\Z)$ as a}%
bilateral sequences $x = (\dots,x_{-2},x_{-1},x_0,x_1,x_2,\dots)$ with the obvious inner product.
\end{exam}

\begin{exam}\label{000150014} For $a < b$ let $\fml C([a,b])$ be the family of all
 \index{C@$\fml C([a,b])$!as an inner product space}%
 \index{inner product!space!$\fml C([a,b])$ as a}%
continuous complex valued functions on the interval $[a,b]$.  For every $f$, $g \in \fml C([a,b])$ define
   \[ \langle f,g \rangle = \int_a^b f(x) \conj{g(x)}\,dx. \]
Then $\fml C([a,b])$ is an inner product space.
\end{exam}

Here is the single most useful fact about semi-inner products. It is an inequality attributed variously to Cauchy,
Schwarz, Bunyakowsky, or combinations of the three.

\begin{thm}[Schwarz inequality]\label{00015002} Let $s$ be a semi-inner product defined on a vector space~$V$.
 \index{Schwarz inequality}%
 \index{inequality!Schwarz}%
Then the inequality
   \[ \abs{s(x,y)} \le \norm x \,\norm y. \]
holds for all vectors $x$ and~$y$ in~$V$.
\end{thm}

\begin{proof}[\emph{Hint for proof}] Fix $x$, $y \in V$.  For every $\alpha \in \K$ we know that
  \begin{equation}\label{Schwarz_ineq}
      0 \le s(x - \alpha y, x - \alpha y) .
  \end{equation}
Expand the right hand side of~\eqref{Schwarz_ineq} into four terms and write $s(y,x)$ in polar form: $s(y,x) = r e^{i\theta}$,
where $r > 0$ and $\theta \in \R$.  Then in the resulting inequality consider those $\alpha$ of the form $e^{-i\theta}t$ where
$t \in \R$.  Notice that now the right side of~\eqref{Schwarz_ineq} is a quadratic polynomial in~$t$.  What can you say about its
discriminant?  \ns
\end{proof}

\begin{prop} Let $V$ be a vector space with semi-inner product $s$ and let $z \in V$.  Then $s(z,z) = 0$
if and only if $s(z,y) = 0$ for all $y \in V$.
\end{prop}

\begin{prop} Let $V$ be a vector space on which a semi-inner product $s$ has been defined.  Then the set
$L := \{z \in V\colon s(z,z) = 0\}$ is a vector subspace of $V$ and the quotient vector space $V/L$ can be
made into an inner product space by defining
    \[ \langle\, [x],[y]\, \rangle := s(x,y) \]
for all $[x]$, $[y] \in V/L$.
\end{prop}

The next proposition shows that an inner product is continuous in its first variable.  Conjugate symmetry then guarantees
continuity in the second variable.

\begin{prop} If $(x_n)$ is a sequence in an inner product space $V$ which converges to a
vector $a \in V$, then $\langle x_n,y \rangle \sto \langle a,y \rangle$ for every  $y \in V$.
\end{prop}

\begin{exam} If $(a_k)$ is a square summable sequence of real numbers, then the series $\sum_{k=1}^\infty k^{-1}a_k$
converges absolutely.
\end{exam}

\begin{exer} Let $0 < a < b$.
 \begin{enumerate}
  \item If $f$ and $g$ are continuous real valued functions on the
interval $[a,b]$, then
   \[\left(\int_a^b f(x)g(x)\,dx\right)^2
             \le \int_a^b\bigl(f(x)\bigr)^2\,dx \cdot
                       \int_a^b\bigl(g(x)\bigr)^2\,dx.\]

\vskip 3 pt

  \item Use part (a) to find numbers $M$ and $N$ (depending on $a$
and~$b$) such that
   \[M(b-a) \le \ln\left(\frac ba\right) \le N(b-a).\]

\vskip 3 pt

  \item Use part (b) to show that $2.5 \le e \le 3$.
 \end{enumerate}
\end{exer}

\begin{proof}[\emph{Hint for proof}] For (b) try $f(x) = \sqrt x$ and $g(x) = \dfrac1{\sqrt x}$.  Then try
$f(x) = \dfrac1x$ and $g(x) = 1$.  For (c) try $a = 1$ and $b = 3$.  Then try $a = 2$ and $b = 5$.  \ns
\end{proof}

\begin{defn} Let $V$ be a vector space. A function
 \index{<unopn@$\norm x$ (norm of a vector~$x$)}%
$\norm{\hphantom{x}} \colon V \sto \R \colon x \mapsto \norm{x}$ is a
 \index{norm}%
\df{norm} on $V$ if
 \begin{enumerate}
  \item $\norm{x + y} \le \norm{x} + \norm{y}$ \qquad  for all $x$, $y \in V$;
  \item $\norm{\alpha x} = \abs{\alpha}\,\norm{x}$\qquad for all $x \in V$ and
$\alpha \in \K$; and
  \item if $\norm{x} = 0$, then $x = \vc 0$.
 \end{enumerate}
The expression $\norm{x}$ may be read as ``the \emph{norm} of $x$'' or ``the
 \index{length}%
\emph{length} of $x$''.  If the function $\norm{\hphantom{x}}$ satisfies~(a) and~(b)
above (but perhaps not~(c)) it is a
 \index{seminorm}%
\df{seminorm} on~$V$.

A vector space on which a norm has been defined is a
 \index{normed!linear space}%
 \index{normed!vector space|seeonly{normed linear space}}%
 \index{space!normed linear}%
\df{normed linear space} (or
 \index{vector!space!normed}%
\df{normed vector space}). A vector in a normed linear space which has norm $1$ is a
 \index{unit!vector}%
 \index{vector!unit}%
\df{unit vector}.
\end{defn}

\begin{exer} Show why it is clear from the definition that $\norm{\vc 0} = 0$ and that norms
(and seminorms) can take on only positive values.
\end{exer}

Every inner product space is a normed linear space.  In more precise language, an inner product on a vector space induces
 \index{norm!induced by an inner product}%
 \index{induced!norm}
a norm on the space.

\begin{prop}\label{00015025} Let $V$ be an inner product space. The map $x \mapsto \norm
x$ defined on $V$ in~\ref{norm_notn} is a norm on~$V$.
\end{prop}

And every normed linear space is a metric space.  Recall the following definition.

\begin{defn}\label{C019614} Let $M$ be a nonempty set. A function $d \colon M \times M \sto \R$ is a
 \index{metric}%
\df{metric} (or
 \index{distance!function}%
 \index{d@$d(x,y)$ (distance between points $x$ and $y$)}%
\df{distance function}) on $M$ if for all $x$, $y$, $z \in M$
 \begin{enumerate}
     \item $d(x,y) = d(y,x)$,
     \item $d(x,y) \le d(x,z) + d(z,y)$, \qquad (the
 \index{triangle inequality}%
 \index{inequality!triangle}%
\emph{triangle inequality})
     \item $d(x,x) = 0$, and
     \item $d(x,y) = 0$ only if $x = y$.
 \end{enumerate}
If $d$ is a metric on a set $M$, then we say that the pair $(M,d)$ is a
 \index{metric!space}%
 \index{space!metric}%
\df{metric space}.  The notation ``$d(x,y)$'' is read, ``the distance between $x$
and~$y$''. (As usual we substitute the phrase, ``Let $M$ be a metric space \dots '' for
the correct formulation ``Let $(M,d)$ be a metric space \dots''.

If $d \colon M \times M \sto \R$ satisfies conditions (1)--(3), but not (4), then $d$ is a
 \index{pseudometric}%
\df{pseudometric} on~$M$.
\end{defn}

As mentioned above, every normed linear space is a metric space. More precisely, a norm on a vector space
induces a
 \index{metric!induced by a norm}%
metric $d$, which is defined by $d(x,y) = \norm{x - y}$.  That is, the distance between
two vectors is the length of their difference.
  \[ \xy
       \Atriangle/<-`<-`>/[``;x`x-y`y]
  \endxy \]
If no other metric is specified we always regard a normed linear space as a metric space
 \index{induced!metric}%
under this induced metric.  The metric is referred to as the \emph{metric induced by the norm}.

\begin{exam}\label{0001503}  Let $V$ be a normed linear space.  Define
$d\colon V \times V \sto \R$ by $d(x,y) = \norm{x - y}$.  Then $d$ is a metric on~$V$.
If $V$ is only a seminormed space, then $d$ is a pseudometric.
\end{exam}

\begin{defn} Vectors $x$ and $y$ in an inner product space $V$ are
 \index{orthogonal}%
\df{orthogonal} (or
 \index{perpendicular}%
\df{perpendicular}) if $\langle x,y \rangle = 0$.  In this case we write
 \index{<binrelperp@$x \perp y$ (orthogonal vectors)}%
$x \perp y$.  Subsets $A$ and $B$ of $V$ are \df{orthogonal} if $a \perp b$ for every $a
\in A$ and $b \in B$.  In this case we write $A \perp B$.
\end{defn}

\begin{defn} If $M$ and $N$ are subspaces of an inner product space $V$ we use the notation
$V = M \oplus N$ to indicate not only that $V$ is the (vector space) direct sum of $M$
and $N$ but also that $M$ and $N$ are orthogonal. Thus we say that $V$ is the
 \index{internal!orthogonal direct sum}%
 \index{direct!sum!orthogonal}%
 \index{orthogonal!direct sum!of inner product spaces}%
(\df{internal}) \df{orthogonal direct sum} of $M$ and~$N$.
\end{defn}

\begin{prop} Let $a$ be a vector in an inner product space $V$.  Then $a \perp x$ for every
$x \in V$ if and only if $a = 0$.
\end{prop}

\begin{prop} In an inner product space $x \perp y$ if and only if $\norm{x + \alpha y} = \norm{x - \alpha y}$
for all scalars $\alpha$.
\end{prop}

\begin{prop}[The Pythagorean theorem]\label{00015037}  If $x \perp y$ in an inner
 \index{Pythagorean theorem}%
product space, then
   \[ \norm{x + y}^2 = \norm x^2 + \norm y^2\,. \]
\end{prop}

\begin{defn}\label{0001504} Let $V$ and $W$ be inner product spaces.  For $(v,w)$ and
$(v',w')$ in $V \times W$ and $\alpha \in \C$ define
 \begin{align*}
      (v,w) + (v',w') &= (v + v', w + w') \\
\intertext{and}
         \alpha(v,w) &= (\alpha v,\alpha w)\,.
 \end{align*}
This results in a vector space, which is the \emph{(external) direct sum} of $V$ and~$W$.
To make it into an inner product space define
   \[ \langle (v,w),(v',w') \rangle = \langle v,v' \rangle + \langle w,w' \rangle. \]
This makes the direct sum of $V$ and $W$ into an inner product space.  It is the
(\df{external orthogonal})
 \index{external!direct sum}%
 \index{direct!sum!external}%
 \index{direct!sum!of inner product spaces}%
 \index{orthogonal!direct sum!of inner product spaces}%
 \index{<binopdirectsum@$\oplus$ (orthogonal direct sum)}%
\df{direct sum} of $V$ and $W$ and is denoted by $V \oplus W$.
\end{defn}

Notice that the same notation $\oplus$ is used for both internal and external direct sums
and for both vector space direct sums (see definition~\ref{000082}) and orthogonal direct
sums. So when we see the symbol $V \oplus W$ it is important to know which category we
are in: vector spaces or inner product spaces, especially as it is common practice to
omit the word ``orthogonal'' as a modifier to ``direct sum'' even in cases when it is
intended.

\begin{exam} In $\R^2$ let $M$ be the $x$-axis and $L$ be the line whose equation is $y = x$.
If we think of $\R^2$ as a (real) vector space, then it is correct to write $\R^2 = M
\oplus L$.  If, on the other hand, we regard $\R^2$ as a (real) inner product space, then
$\R^2 \ne M \oplus L$ (because $M$ and $L$ are not perpendicular).
\end{exam}

\begin{prop}  Let $V$ be an inner product space.  The inner product on $V$, regarded as a map
from $V \oplus V$ into $\C$, is continuous. So is the norm, regarded as a map from $V$
into $\R$.
\end{prop}

Concerning the proof of the preceding proposition, notice that the maps $(v,v') \mapsto
\norm v + \norm{v'}$, $(v,v') \mapsto \sqrt{\norm v^2 + \norm{v'}^2}$, and $(v,v')
\mapsto \max\{\norm v, \norm{v'}\}$ are all norms on $V \oplus V$. Which one is induced
by the inner product on $V \oplus V$?  Why does it not matter which one we use in proving
that the inner product is continuous?

\begin{prop}[The parallelogram law]\label{parallelogram_law}  If $x$ and $y$ are vectors in an inner product
 \index{parallelogram law}%
space, then
   \[ \norm{x + y}^2 + \norm{x - y}^2 = 2\norm x^2 + 2\norm y^2\,. \]
\end{prop}

While every inner product induces a norm not every norm comes from an inner product.

\begin{exam}\label{00015061} There is no inner product on space $\fml C([0,1])$ which induces the uniform norm.
\end{exam}

\begin{proof}[\emph{Hint for proof}] Use the preceding proposition. \ns
\end{proof}

\begin{prop}[The polarization identity]\label{0001507} If $x$ and $y$ are vectors
 \index{polarization identity}%
in a complex inner product space, then
 \[ \langle x,y \rangle = \tfrac14 (\norm{x + y}^2 - \norm{x - y}^2
                       + i\,\norm{x + iy}^2 - i\,\norm{x - iy}^2)\,. \]
\end{prop}

What is the correct identity for a \emph{real} inner product space?

\begin{notn}\label{0001511} Let $V$ be an inner product space, $x \in V$, and $A$,
$B \subseteq V$. If $x \perp a$ for every $a \in A$, we write $x \perp A$; and if $a
\perp b$ for every $a \in A$ and $b \in B$, we write $A \perp B$.
We define $A^\perp$, the
 \index{orthogonal!complement}%
 \index{complement!orthogonal}%
 \index{<unopur@$A^\perp$ (orthogonal complement)}%
\df{orthogonal complement} of $A$, to be $\{x \in V\colon x \perp A\}$.  Since the phrase ``orthogonal complement of $A$''
is a bit unwieldy, especially when used repeatedly, the symbol ``$A^\perp$'' is usually pronounced ``A perp''.
We write $A^{\perp\perp}$ for $\left(A^\perp\right)^\perp$.
\end{notn}

\begin{prop}\label{0001512} If $A$ is a subset of an inner product space $V$, then
$A^\perp$ is a closed linear subspace of~$V$.
\end{prop}

\begin{prop}[Gram-Schmidt orthonormalization] If $(v^k)$ is a linearly independent sequence in
 \index{Gram-Schmidt orthonormalization}%
 \index{orthonormalization}%
an inner product space~$V$, then there exists an orthonormal sequence $(e^k)$ in~$V$ such
that $\spn\{v^1, \dots, v^n\} = \spn\{e^1, \dots, e^n\}$ for every $n \in \N$.
\end{prop}

\begin{cor}\label{0001514} If $M$ is a subspace of a finite dimensional inner product
space $V$, then $V = M \oplus M^\perp$.
\end{cor}

It will be convenient to say of sequences that they \emph{eventually} have a certain property or
that they \emph{frequently} have that property.

\begin{defn}\label{C013121}  Let $(x_n)$ be a sequence of elements of a set $S$ and $P$ be some
property that members of $S$ may possess.  We say that the sequence $(x_n)$
 \index{eventually}%
\df{eventually} has property $P$ if there exists $n_0 \in \N$ such that $x_n$ has
property $P$ for every $n \ge n_0$.  (Another way to say the same thing: $x_n$ has
property $P$ for all but finitely many~$n$.)
\end{defn}

\begin{exam}\label{exam13122} Denote by
 \index{l@$l_c$!sequences which are eventually zero}%
 \index{l@$l_c$!as a vector space}%
$l_c$ the vector space consisting of all sequences of complex numbers which are eventually zero;
that is, the set of all sequences with only finitely many nonzero entries.  They are also referred to as the sequences with
 \index{finite!support}%
 \index{support!finite}%
\emph{finite support}.  The vector space operations are defined pointwise.  We make the space $l_c$ into an
 \index{l@$l_c$!as an inner product space}%
 \index{inner product!space!$l_c$ as a}%
inner product space by
defining $\langle a,b \rangle = \langle\, (a_n),(b_n)\, \rangle := \sum_{k=1}^\infty a_n \conj{b_n}$.
\end{exam}

\begin{defn}\label{C013131} Let $(x_n)$ be a sequence of elements of a set $S$ and $P$ be some
property that members of $S$ may possess.  We say that the sequence $(x_n)$
 \index{frequently}%
\df{frequently} has property $P$ if for every $k \in \N$ there exists $n \ge k$ such that
$x_n$ has property~$P$. (An equivalent formulation: $x_n$ has property $P$ for infinitely
many~$n$.)
\end{defn}

\begin{exam}\label{exam13132} Let $V = l_2$ be the inner product space of all square-summable sequences of complex numbers
(see example~\ref{000150013}) and $M = l_c$ (see example~\ref{exam13122}).  Then the conclusion of~\ref{0001514} fails.
\end{exam}

\begin{defn}\label{def_alg_dual} A
 \index{linear!functional}%
 \index{functional!linear}%
\df{linear functional} on a vector space $V$ is a linear map from $V$ into its scalar
field. The set of all linear functionals on $V$ is the
 \index{dual!algebraic}%
 \index{algebraic!dual space}%
 \index{space!algebraic dual}%
(\df{algebraic}) \df{dual space} of~$V$.  We will use the
 \index{<unopur@$V^\#$ (algebraic dual space)}%
 \index{V@$V^\#$ (algebraic dual space)}%
notation $V^{\#}$ for the algebraic dual space.
\end{defn}

\begin{thm}[Riesz-Fr\'echet Theorem]\label{0001601} If $f \in V^{\#}$ where $V$ is a
 \index{Riesz-Fr\'echet theorem!finite dimensional version}%
finite dimensional inner product space, then there exists a unique vector $a$ in $V$ such
that
    \[ f(x) = \langle x,a \rangle \]
for all $x$ in $V$.
\end{thm}

We will prove shortly (in~\ref{000342})that every continuous linear functional on an \emph{arbitrary} inner
product space has the above representation. The finite dimensional version stated here is a special case, since
every linear map on a finite dimensional inner product space is continuous (see proposition~\ref{00015033}).

\begin{defn}\label{0001602} Let $T\colon V \sto W$ be a linear transformation between complex
inner product spaces. If there exists a function $T^*\colon W \sto V$ which satisfies
   \[ \langle T\vc v,\vc w \rangle = \langle \vc v,T^*\vc w \rangle \]
for all $\vc v \in V$ and $\vc w \in W$, then $T^*$ is the
 \index{adjoint!of an operator on an inner product space}%
 \index{<unopurstar@$T^*$ (inner product space adjoint)}%
 \index{T@$T^*$ (inner product space adjoint)}%
\df{adjoint} (or \df{conjugate transpose}, or \df{Hermitian conjugate}) of~$T$.
\end{defn}

\begin{prop}\label{00016021} If $T\colon V \sto W$ is a linear map between finite
dimensional inner product spaces, then $T^*$ exists.
\end{prop}

\begin{proof}[\emph{Hint for proof}] The functional $\phi\colon V \times W \sto \C \colon
(v,w) \mapsto \langle Tv,w \rangle$ is sesquilinear. Fix $w \in W$ and define $\phi_w
\colon V \sto \C \colon v \mapsto \phi(v,w)$. Then $\phi_w \in V^{\#}$. Use the
\emph{Riesz-Fr\'echet theorem}~(\ref{0001601}).  \ns
\end{proof}

\begin{prop} If $T\colon V \sto W$ is a linear map between finite dimensional inner product
spaces, then the function $T^*$ defined above is linear and $T^{**} = T$.
\end{prop}

\begin{thm}[The fundamental theorem of linear algebra]\label{00016025} If
 \index{fundamental!theorem of linear algebra}%
$T\colon V \sto W$ is a linear map between finite dimensional inner product spaces, then
   \[ \ker T^* = (\ran T)^\perp \quad \text{ and } \quad \ran T^* = (\ker T)^\perp\,. \]
\end{thm}

\begin{defn}\label{000161} An operator $U$ on an inner product space is
 \index{unitary!operator}%
 \index{operator!unitary}%
\df{unitary} if $UU^* = U^*U = I$, that is if $U^* = U^{-1}$.
\end{defn}

\begin{defn}\label{0001612} Two operators $R$ and $T$ on an inner product space $V$ are
 \index{unitary!equivalence}%
 \index{equivalent!unitarily}%
\df{unitarily equivalent} if there exists a unitary operator $U$ on $V$ such that $R =
U^*TU$.
\end{defn}

\begin{prop} If $V$ is an inner product space, then unitary equivalence is in fact an
equivalence relation on~$\ofml L(V)$.
\end{prop}

\begin{defn} An operator $T$ on a finite dimensional inner product space $V$ is
 \index{unitary!diagonalization}%
 \index{diagonalizable!unitarily}%
\df{unitarily diagonalizable} if there exists an orthonormal basis for $V$ with respect
to which $T$ is diagonal. Equivalently, an operator on a finite dimensional inner product
space \emph{with basis} is diagonalizable if it is unitarily equivalent to a diagonal
operator.
\end{defn}

\begin{defn} An operator $T$ on an inner product space is
 \index{self-adjoint!operator}%
 \index{operator!self-adjoint}%
\df{self-adjoint} (or
 \index{Hermitian!operator}%
 \index{operator!Hermitian}%
\df{Hermitian}) if $T^* = T$.
\end{defn}

\begin{defn}\label{000164} A projection $P$ in an inner product space is an
 \index{projection!orthogonal}%
 \index{orthogonal!projection!in an inner product space}%
 \index{operator!orthogonal projection}%
\df{orthogonal projection} if it is self-adjoint.  If $M$ is the range of an orthogonal
projection we will adopt the notation
 \index{PM@$\sbsb PM$ (orthogonal projection onto~$M$)}%
$\sbsb PM$ for the projection rather than the more cumbersome~$\sbsb E{M^\perp M}$.
\end{defn}

\begin{cau} A projection on a vector space or a normed linear space is linear and idempotent,
while an orthogonal projection on an inner product space is linear, idempotent, and
self-adjoint.  This otherwise straightforward situation is somewhat complicated by a
common tendency to refer to orthogonal projections simply as ``projections''.  In fact,
later in these notes we will adopt this very convention.  In inner product spaces
$\oplus$ usually indicates orthogonal direct sum and ``projection'' usually means
``orthogonal projection''. In many elementary linear algebra texts, where everything
happens in $\R^n$, it can be quite exasperating trying to divine whether on any
particular page the author is treating $\R^n$ as a vector space or as an inner product
space.
\end{cau}

\begin{prop} If $P$ is an orthogonal projection on an inner product space $V$, then we have
the orthogonal direct sum decomposition $V = \ker P \oplus \ran P$.
\end{prop}

\begin{defn} If $I_V = P_1 + P_2 + \dots + P_n$ is a resolution of the identity in an
inner product space~$V$ and each $P_k$ is an orthogonal projection, then we say that $I =
P_1 + P_2 + \dots + P_n$ is an
 \index{resolution of the identity!orthogonal}%
 \index{orthogonal!resolution of the identity}%
\df{orthogonal resolution of the identity}.
\end{defn}

\begin{prop} If $I_V = P_1 + P_2 + \dots + P_n$ is an orthogonal resolution of the identity
on an inner product space $V$, then $V = \bigoplus_{k=1}^n \ran P_k$.
\end{prop}

\begin{defn} An operator $N$ on an inner product space is
 \index{normal!operator}%
 \index{operator!normal}%
\df{normal} if $NN^* = N^*N$.
\end{defn}

Two great triumphs of linear algebra are the \emph{spectral theorem} for operators on a
(complex) finite dimensional inner product space (see \ref{00017}), which gives a simply
stated necessary and sufficient condition for an operator to be unitarily diagonalizable,
and theorem~\ref{00018}, which gives a complete classification of those operators.

\begin{thm}[Spectral Theorem for Finite Dimensional Complex Inner Product Spaces]\label{00017}
Let $T$ be an operator on a finite dimensional inner product space~$V$ with (distinct)
 \index{spectral!theorem!for normal operators on inner product spaces}%
eigenvalues $\lambda_1, \dots, \lambda_n$.  Then $T$ is unitarily diagonalizable if and
only if it is normal.  If $T$ is normal, then there exists an orthogonal resolution of
the identity $I_V = P_1 + \dots + P_n$, where for each $k$ the range of the orthogonal
projection $P_k$ is the eigenspace associated with~$\lambda_k$, and furthermore
     \[ T = \lambda_1P_1 +\dots + \lambda_nP_n\,. \]
\end{thm}

\begin{proof}  See \cite{Roman:2005}, page 227. \ns
\end{proof}

\begin{thm}\label{00018} Two normal operators on a finite dimensional inner product space
 \index{normal!operator!conditions for unitary equivalence}%
are unitarily equivalent if and only if they have the same eigenvalues each with the same
multiplicity; that is, if and only if they have the same characteristic polynomial.
\end{thm}

\begin{proof} See \cite{HoffmanK:1971}, page 357.   \ns
\end{proof}

Much of the remainder of these notes is devoted to the definitely nontrivial adventure of finding
appropriate generalizations of the preceding two results to the infinite dimensional
setting and to the astonishing mathematical landscapes which come into view along the way.

\begin{exer} Let $N = \dfrac13\begin{bmatrix}
                                4+2i & 1-i  & 1-i  \\
                                1-i  & 4+2i & 1-i  \\
                                1-i  & 1-i  & 4+2i \end{bmatrix}$.

\vskip 3 pt
 \begin{enumerate}
  \item The matrix $N$ is normal.

\vskip 3 pt

  \item Thus according to the \emph{spectral theorem} $N$
can be written as a linear combination of orthogonal projections.
Explain clearly how to do this (and carry out the computation).
 \end{enumerate}
\end{exer}














\endinput
